% sec-04: the clinical vehicle, medical oncology lead.
\section{The Clinical Vehicle}

The first four years of the programme run through one study: a Phase 1,
first-in-human, open-label, single-arm 3+3 dose escalation of perioperative
daraxonrasib in up to 18 participants with resectable or borderline-resectable
KRAS G12-mutated PDAC, with a staged eight-arm robotic pancreaticoduodenectomy
and an on-premises language model confined to advisory output
\cite{onpremwhippl,phase1ind}. The operation matters to the pharmacology because
it is the only setting that yields paired human tumour tissue under controlled
timing.

\begin{apptable}
\begin{tabularx}{\textwidth}{>{\raggedright\arraybackslash}p{3.3cm} >{\raggedright\arraybackslash}p{3.2cm} >{\raggedright\arraybackslash}X}
\headrow \thead{Endpoint} & \thead{Class} & \thead{Read at}\\
\midrule
Dose-limiting toxicity, MTD or RP2D & Co-primary & Each 3+3 cohort, 28-day window \\
Device- or procedure-related SAEs & Co-primary & 30 days after resection \\
Plasma pharmacokinetics & Secondary & Days $-14$, $-1$, and $+1$ \\
Tumour pathway pharmacodynamics & Secondary & Resection specimen, day 0 \\
ctDNA clearance and pathologic response & Exploratory & Days $+30$ and $+90$; specimen at day 0 \\
Overall survival & Exploratory & 24 months \\
\end{tabularx}
\end{apptable}
\tabcap{The endpoint set, ordered by class rather than by interest. The two
co-primary endpoints govern the dose decision; every pharmacologic endpoint that
motivates the programme is secondary or exploratory.}

\vspace{0.1em}
\noindent\footnotesize\textbf{Positioning note.} \textit{First in human refers
to the integrated surgical and advisory workflow, not to daraxonrasib, which is
investigational and already in Phase 3 evaluation. The robotic configuration is
specified by manufacturer, model, cleared configuration, arm count, and software
version at the site agreement, and no drug supply or letter of authorization is
in place with the agent's developer.}\normalsize\par
